\documentclass[10pt,a4paper]{amsart}
\usepackage{amsmath,amssymb,amsthm,graphicx}
% Retained preamble for the Stacks Project chapter morphisms.tex (Morphisms of Schemes).
% The chapter file pinned-original-source.tex carries only \input{preamble} on its first line, so no
% preamble text lives inside the chapter itself.  The Stacks Project uses one single shared preamble for
% every chapter; the snapshot stages it verbatim in sources/stacks-chow/preamble.tex.  The block below is
% copied mechanically, line for line, from that file's lines 14-18 (the Xy-pic setup) and lines 154-217
% (the native theorem declarations together with the mathematical macros).  Nothing was added, removed or
% redefined.  The wrapper class and the amsmath/amssymb/amsthm/graphicx base packages are supplied by the
% assembler, so the class/\externaldocument part of the shared preamble is not repeated here.
% For commutative diagrams we use Xy-pic
\usepackage[all]{xy}

% We use 2cell for 2-commutative diagrams.
\xyoption{2cell}
% Theorem environments.
%
\theoremstyle{plain}
\newtheorem{theorem}[subsection]{Theorem}
\newtheorem{proposition}[subsection]{Proposition}
\newtheorem{lemma}[subsection]{Lemma}

\theoremstyle{definition}
\newtheorem{definition}[subsection]{Definition}
\newtheorem{example}[subsection]{Example}
\newtheorem{exercise}[subsection]{Exercise}
\newtheorem{situation}[subsection]{Situation}

\theoremstyle{remark}
\newtheorem{remark}[subsection]{Remark}
\newtheorem{remarks}[subsection]{Remarks}

\numberwithin{equation}{subsection}

% Macros
%
\def\lim{\mathop{\mathrm{lim}}\nolimits}
\def\colim{\mathop{\mathrm{colim}}\nolimits}
\def\Spec{\mathop{\mathrm{Spec}}}
\def\Hom{\mathop{\mathrm{Hom}}\nolimits}
\def\Ext{\mathop{\mathrm{Ext}}\nolimits}
\def\SheafHom{\mathop{\mathcal{H}\!\mathit{om}}\nolimits}
\def\SheafExt{\mathop{\mathcal{E}\!\mathit{xt}}\nolimits}
\def\Sch{\mathit{Sch}}
\def\Mor{\mathop{\mathrm{Mor}}\nolimits}
\def\Ob{\mathop{\mathrm{Ob}}\nolimits}
\def\Sh{\mathop{\mathit{Sh}}\nolimits}
\def\NL{\mathop{N\!L}\nolimits}
\def\CH{\mathop{\mathrm{CH}}\nolimits}
\def\proetale{{pro\text{-}\acute{e}tale}}
\def\etale{{\acute{e}tale}}
\def\QCoh{\mathit{QCoh}}
\def\Ker{\mathop{\mathrm{Ker}}}
\def\Im{\mathop{\mathrm{Im}}}
\def\Coker{\mathop{\mathrm{Coker}}}
\def\Coim{\mathop{\mathrm{Coim}}}

% Boxtimes
%
\DeclareMathSymbol{\boxtimes}{\mathbin}{AMSa}{"02}

%
% Macros for moduli stacks/spaces
%
\def\QCohstack{\mathcal{QC}\!\mathit{oh}}
\def\Cohstack{\mathcal{C}\!\mathit{oh}}
\def\Spacesstack{\mathcal{S}\!\mathit{paces}}
\def\Quotfunctor{\mathrm{Quot}}
\def\Hilbfunctor{\mathrm{Hilb}}
\def\Curvesstack{\mathcal{C}\!\mathit{urves}}
\def\Polarizedstack{\mathcal{P}\!\mathit{olarized}}
\def\Complexesstack{\mathcal{C}\!\mathit{omplexes}}
% \Pic is the operator that assigns to X its picard group, usage \Pic(X)
% \Picardstack_{X/B} denotes the Picard stack of X over B
% \Picardfunctor_{X/B} denotes the Picard functor of X over B
\def\Pic{\mathop{\mathrm{Pic}}\nolimits}
\def\Picardstack{\mathcal{P}\!\mathit{ic}}
\def\Picardfunctor{\mathrm{Pic}}
\def\Deformationcategory{\mathcal{D}\!\mathit{ef}}
\title{S-derivations of a module sheaf on an affine scheme correspond to R-derivations of its module of global sections}
\author{The Stacks Project Authors}
\date{stacks-project:morphisms.tex:local-snapshot-20261012; GNU FDL 1.2 or later}
\begin{document}
\maketitle
\noindent\textit{Editorial note.} All delivered mathematical text is the author's own, copied from the pinned Stacks Project chapter file pinned-original-source.tex; no statement or proof step was rewritten, reordered or removed.\par\medskip
\noindent\textit{Editorial note.} Following the multi-claim convention for this source, source\_url is the source-level URL https://stacks.math.columbia.edu/ and the source label lemma-affine-case-derivation is carried in claim\_key and here rather than in a per-claim tag URL.\par\medskip
\noindent\textit{Editorial note.} The preamble is the chapter's shared Stacks Project preamble, retained as sources/stacks-morphisms/retained-preamble.tex: a line-for-line excerpt of the shared Stacks preamble staged with this snapshot. No macro was invented or redefined, and this card adds no preamble text of its own.\par\medskip
\noindent\textit{Editorial note.} Statement span 6408-6417 and proof span 6419-6438 come from the halfA pool card operations/pool-cards/stacks-morphisms-halfA.json and were re-read in the pinned chapter file before this card was written.\par\medskip
\noindent\textit{Editorial note.} The author's proof ends with 'Details omitted.' after displaying the formula defining the derivation. That sentence is kept verbatim in the counted proof span, and what it leaves to the reader is the routine verification that the displayed formula satisfies the derivation properties on standard opens; uniqueness and the explicit construction are both carried out inside the counted span.\par\medskip
\noindent\textit{Allowed standard prerequisites (editorial).}\begin{itemize}
\item The sheaf of relative differentials picture of a derivation: an S-derivation of a sheaf of O\_X-modules F is an R-linear map on sections satisfying the Leibniz rule, and it is determined by its values on a basis of opens. Used at statement lines 6410-6416 and in the uniqueness argument at proof lines 6421-6423. Supplies the objects being classified and the Leibniz rule that determines the values on standard opens.
\item Structure of sections over a standard open: every element of the sections of O\_X over the standard open defined by f is of the form a/f\textasciicircum{}n, and f acts invertibly on the sections of F over that open. Used at proof lines 6420-6424. Makes the displayed Leibniz identity solvable for D(a/f\textasciicircum{}n) and hence makes the derivation determined by D on global sections.
\end{itemize}\medskip
\begin{lemma}
\label{lemma-affine-case-derivation}
Let $R \to A$ be a ring map. Let $\mathcal{F}$
be a sheaf of $\mathcal{O}_X$-modules
on $X = \Spec(A)$. Set $S = \Spec(R)$.
The rule which associates to an $S$-derivation on $\mathcal{F}$
its action on global sections defines a bijection between
the set of $S$-derivations of $\mathcal{F}$ and the set of
$R$-derivations on $M = \Gamma(X, \mathcal{F})$.
\end{lemma}

\begin{proof}
Let $D : A \to M$ be an $R$-derivation. We have to show there exists
a unique $S$-derivation on $\mathcal{F}$ which gives rise to
$D$ on global sections. Let $U = D(f) \subset X$ be a standard affine open.
Any element of $\Gamma(U, \mathcal{O}_X)$ is of the form
$a/f^n$ for some $a \in A$ and $n \geq 0$. By the Leibniz rule
we have
$$
D(a)|_U = a/f^n D(f^n)|_U + f^n D(a/f^n)
$$
in $\Gamma(U, \mathcal{F})$. Since $f$ acts invertibly
on $\Gamma(U, \mathcal{F})$ this completely determines
the value of $D(a/f^n) \in \Gamma(U, \mathcal{F})$.
This proves uniqueness. Existence follows by simply defining
$$
D(a/f^n) := (1/f^n) D(a)|_U - a/f^{2n} D(f^n)|_U
$$
and proving this has all the desired properties (on the basis
of standard opens of $X$). Details omitted.
\end{proof}

\end{document}
